\begin{frame}{Обзор исследований: генерация тестов с LLM}

\textbf{Подход 1: LLM генерирует тесты по требованиям}

\vspace{0.3cm}

\begin{itemize}
    \item \textbf{Alagarsamy et al., 2025 (JSS)}
    \begin{itemize}
        \item Fine-tuned GPT-3.5 генерирует unit-тесты из текста требований
        \item \textit{Проблема:} лишь 67\% тестов точно отражают требования. Проанализируем почему.
    \end{itemize}

    \vspace{0.2cm}

    \item \textbf{Schäfer et al., 2024 (IEEE TSE)}
    \begin{itemize}
        \item TestPilot: LLM генерирует тесты по сигнатуре и коду функции
        \item \textit{Проблема:} сгенерированные тесты слабо совпадают с эталонными ($<$50\%), сложно оценить их смысловую полноту
    \end{itemize}

    \vspace{0.2cm}

    \item \textbf{Ahmed et al., 2024 (arXiv — TSE draft)}
    \begin{itemize}
        \item Auto-TDD: генерация тестов из GitHub-issues
        \item \textit{Проблема:} только 23.6\% тестов реально обнаруживают баг (fail-to-pass), остальные тривиальны. 
    \end{itemize}
\end{itemize}

\end{frame}